\section{Deferred Proofs}
\label{appx:proof}
\label{appx:auxiliary-proofs}

This appendix proves the claims in Sections~\ref{sec:blocks}
and~\ref{sec:together}, the packed PCS claim in
Appendix~\ref{appx:pcs}, and the binary Hadamard claim in
Appendix~\ref{appx:binary-hadamard}, using the notation of each statement.

\subsection{Proof of Proposition~\ref{proposition:ip-complexity}}
\label{appx:proof-ip}

\begin{proof}
\noindent\textbf{Completeness.}
Expand Eq.~\eqref{eq:ip-laurent-polynomial}:
\[
    f(X)
    =
    \sum_{\tau=-(\kappa-1)}^{\kappa-1}
    f_\tau X^\tau,
\]
where:
\[
    f_\tau
    =
    \sum_{\substack{j,j'\in[\kappa]\\j-j'=\tau}}
    \left\langle
        \bm a^{(j)},\bm b^{(j')}
    \right\rangle.
\]
The constant coefficient is the claimed inner product:
\[
    f_0
    =
    \sum_{j=1}^{\kappa}
    \left\langle
        \bm a^{(j)},\bm b^{(j)}
    \right\rangle
    =
    \langle\bm a,\bm b\rangle.
\]
Define the two auxiliary polynomials by:
\[
    h_+(X)
    \coloneqq
    \sum_{\tau=1}^{\kappa-1}
    f_\tau X^{\tau-1},
    \qquad
    h_-(X)
    \coloneqq
    \sum_{\tau=1}^{\kappa-1}
    f_{-\tau}X^{\tau-1}.
\]
They satisfy:
\[
    f(X)
    =
    f_0
    +Xh_+(X)
    +X^{-1}h_-(X^{-1}),
\]
which proves completeness.

\smallskip\noindent\textbf{Soundness.}
Suppose \(c\neq\langle\bm a,\bm b\rangle\). Fix the prover's
polynomials \(h_+\) and \(h_-\) before sampling \(\ext{x}\).
The residual is:
\[
    f(X)-c-Xh_+(X)-X^{-1}h_-(X^{-1}).
\]
Its constant coefficient is \(f_0-c\neq0\). Multiplication by
\(X^{\kappa-1}\) therefore gives a nonzero polynomial of degree
at most \(2\kappa-2\). The evaluation check in
Eq.~\eqref{eq:packed-inner-product-check} accepts with probability
at most:
\[
    \frac{2\kappa-2}{|\ExtensionField|-1}.
\]
By Fact~\ref{fact:ip}, the subsequent length-\(m\)
inner-product reduction has soundness error at most
\(2\log_2m/|\ExtensionField|\). The claimed bound follows from
the union bound.

\smallskip\noindent\textbf{Efficiency.}
Computing \(h_+\) and \(h_-\) requires one length-\(m\)
base-field inner product for each ordered pair \(j\neq j'\).
The number of multiplications of type \(\bb\) is:
\[
    m\kappa(\kappa-1)
    =
    N(\kappa-1).
\]

Each packed vector requires \(m\kappa=N\) mixed-field
multiplications, for a total of \(2N\,\be\).

By Fact~\ref{fact:ip}, the number of multiplications of type
\(\ee\) in the length-\(m\) inner-product reduction is:
\[
    4m+O(\log m).
\]
Computing the positive and negative powers of \(\ext{x}\)
contributes \(O(\kappa)\)
additional \(\ee\) operations. Substituting \(m=N/\kappa\) gives:
\[
    N(\kappa-1)\,\bb
    +
    2N\,\be
    +
    \left(
        \frac{4N}{\kappa}
        +
        O\left(
            \kappa+\log\frac{N}{\kappa}
        \right)
    \right)\ee.
\]

The matrices and packed vectors occupy \(O(N)\) base-field
coordinates. Streaming the cross-block inner products retains
only the \(2\kappa-2\) auxiliary coefficients; the challenge
powers occupy \(O(\kappa^2)\) coordinates. An in-place
inner-product reduction therefore uses \(O(N+\kappa^2)\) peak
memory, or \(O(N)\) when \(m\geq\kappa\).
\end{proof}

\subsection{Proof of Proposition~\ref{proposition:hadamard-complexity}}
\label{appx:proof-hadamard}

\begin{proof}
\noindent\textbf{Completeness.}
Define the correct column claims by:
\[
    \ext{d}_j
    \coloneqq
    \left\langle
        \ext{\bm l},
        \bm a^{(j)}\circ\bm b^{(j)}
    \right\rangle
    \qquad
    \text{for }j\in[\kappa].
\]
An honest prover sends \(\ext{\bm c}=\ext{\bm d}\), so:
\[
    \ext{c}
    =
    \left\langle
        \ext{\bm c},\ext{\bm r}
    \right\rangle
    =
    \ext{\bm l}^{\top}
    (\bm A\circ\bm B)
    \ext{\bm r}.
\]
The residual then vanishes at every interpolation node:
\[
    f(\alpha_j)
    =
    \ext{d}_j-\ext{c}_j
    =
    0
    \qquad
    \text{for every }j\in[\kappa].
\]
Thus \(\Psi(X)\) divides \(f(X)\). The degree bounds are:
\[
    \deg f\leq2\kappa-2
    \qquad\text{and}\qquad
    \deg\Psi=\kappa.
\]
Consequently, the quotient has degree at most \(\kappa-2\):
\[
    h(X)=\frac{f(X)}{\Psi(X)}.
\]

Evaluation at \(\ext{x}\) gives the inner-product claim:
\[
\begin{aligned}
    \ext{d}
    &=
    g_C(\ext{x})
    +
    \Psi(\ext{x})h(\ext{x})\\
    &=
    \left\langle
        \ext{\bm a},
        \ext{\bm b}\circ\ext{\bm l}
    \right\rangle.
\end{aligned}
\]
To verify the terminal claims, write:
\[
    \ext{\bm\chi}_+
    \coloneqq
    \tensorspan{\ext{\bm y}'},
    \qquad
    \ext{\bm\chi}_-
    \coloneqq
    \tensorspan{\ext{\bm y}'}.
\]
The terminal scalars from Fact~\ref{fact:ip} satisfy:
\[
\begin{aligned}
    \ext{u}_1
    &=
    \left\langle
        \ext{\bm a},\ext{\bm\chi}_+
    \right\rangle=
    \sum_{j=1}^{\kappa}
    \psi_j(\ext{x})
    \left\langle
        \bm a^{(j)},\ext{\bm\chi}_+
    \right\rangle
    =
    \left\langle
        \ext{\bm v}_a,\ext{\bm\rho}
    \right\rangle,
\end{aligned}
\]
and:
\[
\begin{aligned}
    \ext{u}_2
    &=
    \left\langle
        \ext{\bm b}\circ\ext{\bm l},
        \ext{\bm\chi}_-
    \right\rangle
    =
    \sum_{j=1}^{\kappa}
    \psi_j(\ext{x})
    \left\langle
        \bm b^{(j)},
        \ext{\bm l}\circ\ext{\bm\chi}_-
    \right\rangle
    =
    \left\langle
        \ext{\bm v}_b,\ext{\bm\rho}
    \right\rangle.
\end{aligned}
\]
These are the two Lagrange consistency checks.

For the column aggregation, write:
\[
    \ext{\bm\chi}''
    \coloneqq
    \tensorspan{\ext{\bm y}''}.
\]
The first aggregated claim is:
\[
\begin{aligned}
    \ext{v}_1
    &=
    \left\langle
        \ext{\bm v}_a,\ext{\bm\chi}''
    \right\rangle
    =
    \left\langle
        \bm a,
        \ext{\bm\chi}''
        \kronecker
        \ext{\bm\chi}_+
    \right\rangle
    =
    \left\langle
        \bm a,
        \ext{\bm\xi}_a
    \right\rangle.
\end{aligned}
\]
For the second claim, multiply the two-entry row factors:
\[
    \ext{\bm l}\circ\ext{\bm\chi}_-
    =
    \bigotimes_{j=1}^{\log_2m}
    \begin{pmatrix}
        \ext l_{j,0}(1-\ext y'_j)\\
        \ext l_{j,1}\ext y'_j
    \end{pmatrix}.
\]
Hence:
\[
\begin{aligned}
    \ext{v}_2
    &=
    \left\langle
        \ext{\bm v}_b,\ext{\bm\chi}''
    \right\rangle
    =
    \left\langle
        \bm b,
        \ext{\bm\chi}''
        \kronecker
        \left(
            \ext{\bm l}\circ\ext{\bm\chi}_-
        \right)
    \right\rangle
    =
    \left\langle
        \bm b,
        \ext{\bm\xi}_b
    \right\rangle,
\end{aligned}
\]
which proves completeness of the terminal tensor-weighted claims.
These are valid packed queries under
Eq.~\eqref{eq:pcs-opening-relation}, including when a factor
cannot be normalized to an MLE weight.

\smallskip\noindent\textbf{Soundness.}
Suppose the weighted Hadamard claim is false and the scalar check
accepts:
\[
    \ext{c}
    =
    \left\langle
        \ext{\bm c},\ext{\bm r}
    \right\rangle.
\]
Then \(\ext{\bm c}\neq\ext{\bm d}\), so \(f(\alpha_j)\neq0\)
for some \(j\in[\kappa]\). For any \(h\) fixed before
\(\ext{x}\), consider the residual:
\[
    f(X)-\Psi(X)h(X).
\]
It is nonzero and has degree at most \(2\kappa-2\).
Consequently, a uniformly random
\(\ext{x}\in\ExtensionField\setminus\mathcal H\) satisfies the
quotient identity with probability at most:
\[
    \frac{2\kappa-2}
         {|\ExtensionField|-\kappa}.
\]

By Fact~\ref{fact:ip}, the subsequent length-\(m\)
inner-product reduction contributes soundness error at most:
\[
    \frac{2\log_2m}
         {|\ExtensionField|}.
\]

The last step authenticates the column values sent after the
inner-product reduction. Their correct values are:
\[
\begin{aligned}
    \widehat v_{a,j}
    &\coloneqq
    \left\langle
        \bm a^{(j)},\ext{\bm\chi}_+
    \right\rangle,
    \qquad
    \widehat v_{b,j}
    \coloneqq
    \left\langle
        \bm b^{(j)},
        \ext{\bm l}\circ\ext{\bm\chi}_-
    \right\rangle.
\end{aligned}
\]
If either
\(\ext{\bm v}_a\neq\widehat{\bm v}_a\) or
\(\ext{\bm v}_b\neq\widehat{\bm v}_b\), the corresponding
difference vector defines a nonzero multilinear polynomial in
\(\log_2\kappa\) variables. Since these vectors are fixed before
\(\ext{\bm y}''\) is sampled, each incorrect terminal claim
passes the randomized MLE aggregation with probability at most:
\[
    \frac{\log_2\kappa}{|\ExtensionField|}.
\]
A union bound over the two claims gives the additional error:
\[
    \frac{2\log_2\kappa}{|\ExtensionField|}.
\]
Combining the three error terms proves the stated soundness bound.
The same argument holds for \(\mathcal H\subseteq\mathbb K
\subseteq\ExtensionField\): distinct nodes, polynomial degrees,
and challenge order are unchanged. Lookup tables compute the
same quotient and introduce no additional soundness error.

\smallskip\noindent\textbf{Efficiency.}
The counts below concern \(\mathcal H\subseteq\Field\).
Appendix~\ref{appx:hadamard-comparison} gives the untabulated
and tabulated intermediate-field costs.
Computing
\(\ext{\bm c}\) requires \(N\) multiplications of type \(\bb\)
to form \(\bm A\circ\bm B\), followed by \(N\) multiplications
of type \(\be\) to contract the result with
\(\ext{\bm l}\).

For each row, interpolating
\(g_{A,i}(X)\) and \(g_{B,i}(X)\) requires
\(2\kappa^2\) multiplications of type \(\bb\), while multiplying
the resulting degree-\((\kappa-1)\) polynomials requires another
\(\kappa^2\) multiplications of type \(\bb\). The total number
of base-field multiplications is:
\[
    N+3m\kappa^2
    =
    (3\kappa+1)N.
\]

Each product
\(g_{A,i}(X)g_{B,i}(X)\) contains \(2\kappa-1\)
coefficients. Contracting them with \(\ext{l}_i\) over all rows
uses the following number of mixed-field multiplications:
\[
    m(2\kappa-1)
    =
    2N-m.
\]
Interpolating \(g_C(X)\) and dividing the residual by \(\Psi(X)\)
require
\(O(\kappa^2)\) additional \(\be\) operations.

After \(\ext{x}\) is sampled, forming
\(\ext{\bm a}=\bm A\ext{\bm\rho}\) and
\(\ext{\bm b}=\bm B\ext{\bm\rho}\) requires \(2N\,\be\).
After the inner-product reduction, computing the two vectors of
column values \(\ext{\bm v}_a\) and \(\ext{\bm v}_b\) requires
another \(2N\,\be\). Including the \(N\,\be\) operations used to
construct \(\ext{\bm c}\), the leading mixed-field multiplication
count is:
\[
\begin{aligned}
    &N+(2N-m)+2N+2N\\
    &\qquad=
    7N-m
    =
    7N-\frac{N}{\kappa}.
\end{aligned}
\]
The remaining mixed-field operations contribute \(O(\kappa^2)\).

Forming
\(\ext{\bm b}\circ\ext{\bm l}\) before the inner-product
reduction requires \(m\,\ee\). Fact~\ref{fact:ip} gives the
following extension-field multiplication count for the reduction:
\[
    4m+O(\log m).
\]
The second column query uses:
\[
    \ext{\bm l}
    \circ
    \tensorspan{\ext{\bm y}'}.
\]
Forming this vector requires another \(m\,\ee\).
Expanding the supplied two-entry factors of \(\ext{\bm l}\)
uses at most \(2(m-1)\,\ee\); generating
\(\tensorspan{\ext{\bm y}'}\) uses at most \((m-1)\,\ee\).
Evaluating the Lagrange basis and the degree-\(O(\kappa)\)
polynomials, and forming the random column aggregations, requires
\(O(\kappa)\) further \(\ee\) operations. The total number
of extension-field multiplications is therefore at most:
\[
    9m+O(\kappa+\log m)
    =
    \frac{9N}{\kappa}
    +
    O\left(
        \kappa+\log\frac{N}{\kappa}
    \right).
\]

The matrices \(\bm A\) and \(\bm B\) occupy \(2N\) base-field
coordinates. Processing the rows sequentially retains
\(O(\kappa)\) extension-field coefficients for the residual and
quotient. The column claims and tensor factors occupy another
\(O(\kappa)\) extension-field elements; the packed vectors and
inner-product reduction use a constant number of length-\(m\)
buffers. The peak memory, in base-field coordinates, is:
\[
    O(N+\kappa^2).
\]
This is \(O(N)\) when \(m\geq\kappa\).
\end{proof}

\subsection{Proof of Proposition~\ref{proposition:binary-hadamard-complexity}}
\label{appx:proof-binary-hadamard}

\begin{proof}
\noindent\textbf{Composition and completeness.}
Choose \(\ext{\bm l}\) and \(\ext{\bm r}\) as independent
random MLE tensor weights of lengths \(m\) and \(\kappa\),
respectively. The prover supplies the projected value
\(\ext c=\ext{\bm l}^{\top}\bm C\ext{\bm r}\), the column
claims \(\ext{\bm c}\), and the mixed-basis quotient coefficients.
The verifier checks \(\ext c=\ext{\bm c}^{\top}\ext{\bm r}\),
then samples \(\ext{\bm x}\in\ExtensionField^s\).
Apply Fact~\ref{fact:ip} to the inner product obtained from
Eq.~\eqref{eq:binary-hadamard-certificate} at \(\ext{\bm x}\).
If \(\ext{\bm y}\in\ExtensionField^{\log_2m}\) is its terminal
point, the three claims to authenticate are:
\[
\begin{aligned}
    \ext c
      &=\ext{\bm l}^{\top}\bm C\ext{\bm r},\\
    \ext v_1
      &=\tensorspan{\ext{\bm y}}^{\top}
          \bm A\bm\psi(\ext{\bm x}),\\
    \ext v_2
      &=\bigl(\ext{\bm l}\circ
          \tensorspan{\ext{\bm y}}\bigr)^{\top}
          \bm B\bm\psi(\ext{\bm x}).
\end{aligned}
\]
Each is a packed opening in the sense of
Eq.~\eqref{eq:pcs-opening-relation}.
In particular, the row query for \(\bm B\) remains a tensor
vector: multiply the corresponding two-entry factors
coordinatewise. Proposition~\ref{proposition:packed-pcs}
therefore authenticates these claims directly, with no additional
column-aggregation step. The product identity and the
completeness of the two invoked protocols give completeness.

\smallskip\noindent\textbf{Prover time.}
The transforms and weighted accumulation, including the column
claims, use \(O(m3^s)\) bit operations and
\(O(m3^s)\) extension-field XORs.
Generating the row and column weights and
\(\bm\psi(\ext{\bm x})\) uses \(O(m+\kappa)\)
extension-field multiplications.
The matrix--vector products with the binary matrices
\(\bm A,\bm B,\bm C\) use \(O(N)\) extension-field XORs.
Computing \(\ext{\bm l}^{\top}(\bm C\ext{\bm r})\),
weighting the packed vector for \(\bm B\), computing the
inner-product value, and applying Fact~\ref{fact:ip} use
\(O(m)\) extension-field multiplications.
The honest prover computes this value from the packed vectors;
only the verifier needs to evaluate the full quotient certificate.
Thus, before the packed openings, the prover uses:
\[
\begin{aligned}
    &O(m3^s)\text{ bit operations},\\
    &O(m3^s)\text{ extension-field XORs},\\
    &O(m+\kappa)\text{ extension-field multiplications}.
\end{aligned}
\]
An extension-field XOR costs \(O(\kappa)\) bit operations,
and a schoolbook extension-field multiplication costs
\(O(\kappa^2)\). Since \(N=m\kappa\) and \(3^s\geq\kappa\),
the preceding bound becomes \(O(N3^s+\kappa^3)\) bit operations.
By Eq.~\eqref{eq:pcs-prover-cost}, each packed-opening reduction
adds \(O(\kappa N+\kappa^3)\) bit operations.
All three are absorbed in the stated prover bound.

\smallskip\noindent\textbf{Verifier time.}
The \(\kappa\) column claims and the \(Q=3^s-\kappa\)
quotient coefficients together form a tensor of \(3^s\)
coefficients in the basis \(\bm\chi\).
Contracting one coordinate at a time evaluates the certificate
in \(O(3^s)\) extension-field operations.
The column check and the inner-product verification require
another \(O(\kappa+\log_2m)\) operations.
The verifier retains the row queries as two-entry tensor factors
and never constructs length-\(m\) tables.
Its cost before openings is therefore
\(O(\kappa^2(3^s+\log_2m))\) bit operations.
Proposition~\ref{proposition:packed-pcs} adds
\(O(\kappa^3(1+\log_2m))\) bit operations for the three
packed-opening reductions, giving the claimed total.

\smallskip\noindent\textbf{Communication and memory.}
The prover sends \(\kappa+Q=3^s\) column and quotient
coefficients, one projected value, and
\(O(1+\log_2m)\) elements for the inner-product reduction.
Each packed opening adds \(O(\kappa+\log_2m)\) elements.
This gives \(O(3^s+\log_2m)\) extension-field elements in total.
For memory, retain the \(O(N)\) input bits, a constant number
of length-\(m\) extension-field buffers, and \(3^s\)
extension-field coefficients. Row-streamed transforms need
only \(O(3^s)\) additional bits.
Each packed-opening reduction uses \(O(N+\kappa^2)\) bits
of workspace, which can be reused across openings.
The peak is \(O(N+\kappa3^s)\) bits.

\smallskip\noindent\textbf{Soundness.}
An accepting execution either contains an accepted false
packed-opening claim or authenticates all three displayed equalities.
For \(\bm A\circ\bm B\neq\bm C\), the initial random
tensor projection of this nonzero binary matrix vanishes with
probability at most \(\log_2N/|\ExtensionField|\).
If this projection is nonzero and the opening for \(\bm C\)
is valid, passing the column check requires a false column claim.
For any quotients fixed before \(\ext{\bm x}\),
the certificate residual is then a nonzero polynomial of total
degree at most \(2s\), giving error at most
\(2s/|\ExtensionField|\).
Fact~\ref{fact:ip} contributes at most
\(2\log_2m/|\ExtensionField|\).
The algebraic errors sum to
\(3\log_2N/|\ExtensionField|\), since
\(\log_2N=s+\log_2m\).
Finally, Eq.~\eqref{eq:pcs-soundness} bounds each packed-opening
error by
\(\epsilon_{\mathsf{PCS}}+(1+2\log_2m)/|\ExtensionField|\).
A union bound over the three openings and the algebraic checks
proves the claimed soundness bound.
\end{proof}

\subsection{Proof of Proposition~\ref{proposition:permutation-complexity}}
\label{appx:proof-permutation}

\begin{proof}
\noindent\textbf{Completeness.}
If \(\bm a\) is a permutation of \(\bm b\), then:
\[
    \prod_{i=1}^{N}(\ext z+a_i)
    =
    \prod_{i=1}^{N}(\ext z+b_i),
\]
for every \(\ext z\in\ExtensionField\). It remains to verify
one grand-product reduction.

The column-product tree satisfies:
\[
    \bm C_{t+1}(\ext z)
    =
    \bm C_{t,L}(\ext z)
    \circ
    \bm C_{t,R}(\ext z),
\]
for every
\(t=0,\ldots,\log_2\kappa-1\). The recursive row-product tree
similarly satisfies:
\[
    \ext{\bm c}_{i+1}
    =
    \ext{\bm c}_{i,L}
    \circ
    \ext{\bm c}_{i,R},
    \qquad
    i=0,\ldots,\log_2m-1,
\]
and hence:
\[
    \ext c_{\log_2m}
    =
    \mathsf{GrandProd}
    \bigl(
        \bm C_{\log_2\kappa}(\ext z)
    \bigr).
\]

Each tree relation reduces to an inner product. Take vectors:
\[
    \ext{\bm u},
    \ext{\bm v},
    \ext{\bm\xi}
    \in\ExtensionField^n.
\]
Write the weighted Hadamard claim as:
\begin{equation}
\begin{aligned}
    \ext c
    &\EqOrNot
    \left\langle
        \ext{\bm\xi},
        \ext{\bm u}\circ\ext{\bm v}
    \right\rangle\\
    &=
    \left\langle
        \ext{\bm u},
        \ext{\bm v}\circ\ext{\bm\xi}
    \right\rangle.
\end{aligned}
\label{eq:permutation-hadamard-as-ip}
\end{equation}
The prover forms
\(
    \ext{\bm v}\circ\ext{\bm\xi}
\)
and applies Fact~\ref{fact:ip} to
Eq.~\eqref{eq:permutation-hadamard-as-ip}.

Recursing through the row tree and then the column tree leaves:
\[
    \left\langle
        \ext{\bm\xi}_0,\bm d
    \right\rangle
    \EqOrNot
    \ext c_0
    -
    \ext z
    \left\langle
        \ext{\bm\xi}_0,\bm 1_N
    \right\rangle,
\]
which is an MLE opening of the original committed vector
\(\bm d\). This proves completeness.

\smallskip\noindent\textbf{Soundness.}
If \(\bm a\) is not a permutation of \(\bm b\), the fingerprint
difference is nonzero:
\[
    Q(Z)
    \coloneqq
    \prod_{i=1}^{N}(Z+a_i)
    -
    \prod_{i=1}^{N}(Z+b_i).
\]
The leading terms cancel, so \(\deg Q\leq N-1\). A uniform
\(\ext z\in\ExtensionField\) is a root with probability at most:
\[
    \min\left\{
        1,
        \frac{N-1}{|\ExtensionField|}
    \right\}.
\]

For one grand-product proof, the row-level-\(i\) relation in
Eq.~\eqref{eq:permutation-hadamard-as-ip} has length:
\[
    \frac{m}{2^{i+1}}.
\]
Fact~\ref{fact:ip} bounds its error by:
\[
    \frac{
        2\log_2(m/2^{i+1})
    }{
        |\ExtensionField|
    }.
\]
Summing over the row tree gives:
\[
\begin{aligned}
    \sum_{i=0}^{\log_2m-1}
    \frac{
        2\log_2(m/2^{i+1})
    }{
        |\ExtensionField|
    }
    &=
    \frac{
        \log_2m(\log_2m-1)
    }{
        |\ExtensionField|
    }.
\end{aligned}
\]

At column level \(t\), the interpolation-domain size is:
\[
    2^{\log_2\kappa-t}.
\]
The corresponding residual degree is at most:
\[
    2^{\log_2\kappa-t+1}-2.
\]
The quotient is fixed before the evaluation challenge. Thus a
false quotient identity passes the polynomial check with probability
at most:
\[
    \frac{
        2^{\log_2\kappa-t+1}-2
    }{
        |\ExtensionField|
        -
        2^{\log_2\kappa-t}
    }.
\]

The evaluated Hadamard claim has length:
\[
    \frac{m}{2^{\log_2\kappa-t}}.
\]
The inner-product reduction of Fact~\ref{fact:ip} therefore
contributes:
\[
    \frac{
        2\log_2
        \left(
            m/2^{\log_2\kappa-t}
        \right)
    }{
        |\ExtensionField|
    }.
\]
The random combination of the resulting column claims
contributes at most:
\[
    \frac{
        2\log_2
        \left(
            2^{\log_2\kappa-t}
        \right)
    }{
        |\ExtensionField|
    }.
\]
These two terms sum to:
\[
    \frac{2\log_2m}{|\ExtensionField|}.
\]
Summing over \(t=0,\ldots,\log_2\kappa-1\) gives
\(\epsilon_{\mathsf{GP}}\). A union bound over the fingerprint,
the two grand-product proofs, and their terminal openings proves
the permutation bound.

\smallskip\noindent\textbf{Efficiency.}
At column level \(t\), constructing the coefficient representation
and quotient requires:
\[
    O\left(
        \frac{
            \kappa^2m
        }{
            2^{\log_2\kappa-t}
        }
    \right)\bb.
\]
Summing gives the following base-field multiplication count:
\[
\begin{aligned}
    \sum_{t=0}^{\log_2\kappa-1}
    O\left(
        \frac{
            \kappa^2m
        }{
            2^{\log_2\kappa-t}
        }
    \right)
    &=
    O\bigl(m\kappa(\kappa-1)\bigr)\\
    &=
    O(N\kappa).
\end{aligned}
\]
Evaluating the coefficient polynomials and constructing the
packed vectors require
\(O(N)\) multiplications of type \(\be\).

For a length-\(n\) weighted Hadamard claim, forming
\(
    \ext{\bm v}\circ\ext{\bm\xi}
\)
requires \(n\) multiplications of type \(\ee\). Fact~\ref{fact:ip}
gives the following count for the inner-product reduction:
\[
    4n+O(\log n).
\]
The full length-\(n\) Hadamard proof therefore uses the following
number of extension-field multiplications:
\[
    5n+O(\log n).
\]

The row-tree inner-product lengths sum to:
\[
    \sum_{i=0}^{\log_2m-1}
    \frac{m}{2^{i+1}}
    =
    m-1.
\]
The row-tree extension-field multiplication count is therefore:
\[
    5(m-1)
    +
    O\bigl((\log_2m)^2\bigr).
\]

The column-tree inner-product lengths sum to:
\[
\begin{aligned}
    \sum_{t=0}^{\log_2\kappa-1}
    \frac{m}{2^{\log_2\kappa-t}}
    &=
    m\frac{\kappa-1}{\kappa}.
\end{aligned}
\]
Their leading extension-field multiplication count is:
\[
    5m\frac{\kappa-1}{\kappa}.
\]
Including challenge generation and polynomial evaluation gives
the total extension-field multiplication count:
\[
\begin{aligned}
    &5(m-1)
    +
    5m\frac{\kappa-1}{\kappa}
    +
    O\left(
        \kappa
        +
        (\log_2m)^2
        +
        \log_2m\log_2\kappa
    \right)\\
    &\qquad=
    \frac{10N}{\kappa}
    +
    O\left(
        \kappa
        +
        (\log_2m)^2
        +
        \log_2m\log_2\kappa
    \right).
\end{aligned}
\]
Processing the coefficient layers sequentially retains the
\(O(N)\) witness coordinates and \(O(\kappa^2)\) coordinates
for quotient coefficients and challenge powers. Peak memory is
\(O(N+\kappa^2)\), or \(O(N)\) when \(m\geq\kappa\).
\end{proof}

\subsection{Proof of Proposition~\ref{proposition:packed-pcs}}
\label{appx:proof-pcs}

\begin{proof}
\noindent\textbf{Binding and completeness.}
Basis packing is injective: each entry of
\(\bm A\ext{\bm\beta}\) uniquely determines its row of
\(\bm A\). Binding therefore follows from the underlying PCS.
Applying a hiding commitment to this fixed encoding also preserves
commitment hiding.

For an honest prover,
\(\bm V=\bm A^\top\bm\Xi\), so the identity in
Eq.~\eqref{eq:pcs-encoded-matrix} and its random projection hold.
It remains to show that the verifier computes the public terminal
value exactly.

By Eq.~\eqref{eq:pcs-multiplication-matrix}, multiplication by
\(\ext z\) acts on a column of base-field coordinates through
\(\bm M_{\ext z}\). This map is \(\Field\)-linear in \(\ext z\),
and \(\bm M_{1}=\bm I_\kappa\).
We prove by backward induction that
\(\bm W_j=\bm U_j^\top\bm\Xi_j\).
The empty suffix has coordinate row \(\bm e_1^\top\), since
\(\ext\beta_1=1\). Hence the initialization is correct:
\[
    \bm U_{\ell+1}^\top\bm\Xi_{\ell+1}
    =\bm e_1\bm e_1^\top
    =\bm W_{\ell+1}.
\]
For the induction step, apply the block identities in
Eq.~\eqref{eq:pcs-query-blocks} and the induction hypothesis:
\[
\begin{aligned}
    \bm U_j^\top\bm\Xi_j
    &=\bm M_{1-\ext y_j}\bm W_{j+1}
        \bm M_{\ext\chi^{(j)}_0}^\top\\
    &\quad{}+\bm M_{\ext y_j}\bm W_{j+1}
        \bm M_{\ext\chi^{(j)}_1}^\top\\
    &=\bm W_{j+1}\bm M_{\ext\chi^{(j)}_0}^\top\\
    &\quad{}+\bm M_{\ext y_j}\bm W_{j+1}
        \bm M_{\ext\chi^{(j)}_1-\ext\chi^{(j)}_0}^\top\\
    &=\bm W_j.
\end{aligned}
\]
The second equality uses the linearity of \(\bm M\), and the
last is Eq.~\eqref{eq:pcs-query-recurrence}.
At \(j=1\), we obtain
\(\bm W=\bm W_1=\bm U^\top\bm\Xi\). Consequently:
\[
    \ext{\bm\beta}^\top\bm W\ext{\bm r}
    =(\bm U\ext{\bm\beta})^\top\bm\Xi\ext{\bm r}
    =\ext{\bm\xi}^\top\bm\Xi\ext{\bm r}.
\]
The scalar check in Eq.~\eqref{eq:pcs-query-matrix} therefore
verifies the query-side terminal claim exactly. Completeness
follows from Fact~\ref{fact:ip} and the underlying PCS opening.
When \(m=1\), the suffix is empty and the initialization alone
gives the required matrix; no recurrence or sum-check round is
needed.

\smallskip\noindent\textbf{Soundness.}
Fix the matrix \(\bm A\) bound by the commitment and suppose that
Eq.~\eqref{eq:pcs-opening-relation} is false.
If the first scalar check accepts, then
\(\ext{\bm v}\neq\bm A^\top\ext{\bm\chi}\), hence
\(\bm V\neq\bm A^\top\bm\Xi\).
By injectivity of basis packing, the following row is nonzero:
\[
    \ext{\bm\beta}^{\top}
    (\bm V-\bm A^\top\bm\Xi)
    \in\ExtensionField^{1\times\kappa}\setminus\{\bm0_\kappa^\top\}.
\]
This row is fixed before \(\ext{\bm r}\) is sampled. Its inner
product with a uniform \(\ext{\bm r}\) vanishes with probability
\(1/|\ExtensionField|\); otherwise,
Eq.~\eqref{eq:pcs-random-opening} is false.

Except with probability \(2\ell/|\ExtensionField|\),
Fact~\ref{fact:ip} then leaves a false terminal claim in
Eq.~\eqref{eq:pcs-terminal-claims}.
The recurrence computes \(\bm W\) from the public tensor factors
and the terminal point, so a false query-side claim is rejected
deterministically. If that check passes, the witness-side claim
must be false. The underlying PCS accepts this adaptively chosen
false MLE claim with probability at most
\(\epsilon_{\mathsf{PCS}}\).
A union bound gives Eq.~\eqref{eq:pcs-soundness}.

\smallskip\noindent\textbf{Prover complexity.}
The prover expands \(\ext{\bm\chi}\) and reads its basis
coordinates to obtain \(\bm\Xi\).
This costs \(O(m)\) extension-field operations.
The matrix product \(\bm V=\bm A^\top\bm\Xi\) takes
\(\kappa N\) base-field multiplications. Its rows are already
the coordinates of the message
\(\ext{\bm v}=\bm V\ext{\bm\beta}\); packing them requires
only a change of representation.

Computing \(\bm\Xi\ext{\bm r}\) takes \(N\) mixed-field
multiplications; the extension-field inner-product reduction takes
\(4m+O(\ell)\) extension-field multiplications.
The target \(\ext c\) can be computed by interpreting each column
of \(\bm V\) as an extension-field element and taking the inner
product of the resulting row with \(\ext{\bm r}\), using
\(\kappa\) extension-field multiplications.
The prover need not construct \(\bm U\) or \(\bm W\).
These costs give Eq.~\eqref{eq:pcs-prover-cost}.
Under schoolbook arithmetic they are
\(O(\kappa N+\kappa^3)\,\bb\), which is
\(O(\kappa N)\,\bb\) when \(m\geq\kappa\).

The matrices \(\bm A,\bm\Xi\) occupy \(O(N)\) base-field
coordinates, and \(\bm V\) and the length-\(\kappa\)
extension-field vectors occupy \(O(\kappa^2)\).
A constant number of length-\(m\) extension-field buffers also
occupy \(O(m\kappa)=O(N)\) base-field coordinates.
Thus peak prover memory, excluding the underlying PCS workspace,
is \(O(N+\kappa^2)\), or \(O(N)\) when \(m\geq\kappa\).

\smallskip\noindent\textbf{Verifier complexity and communication.}
The verifier implements Eq.~\eqref{eq:pcs-query-recurrence}
without forming the multiplication matrices.
Right multiplication by \(\bm M_{\ext z}^\top\) scales each
row, interpreted in the fixed basis as an extension-field element,
by \(\ext z\). Left multiplication by \(\bm M_{\ext z}\)
scales each column in the same way. Each scaling therefore costs
\(\kappa\) extension-field multiplications.
The two terms in each recurrence step require three such scalings,
followed by \(\kappa^2\) base-field additions.
Over \(\ell\) steps, the cost is at most
\(3\kappa\ell\) extension-field multiplications and
\(O(\kappa^2\ell)\) base-field operations, including coordinate
rearrangement. A constant number of intermediate matrices use
\(O(\kappa^2)\) additional base-field coordinates.

The initial scalar check costs \(\kappa\) extension-field
multiplications. For each of \(\ext c\) and the expected terminal
value \(\ext w\), the verifier packs the columns of the relevant
matrix and takes one length-\(\kappa\) extension-field inner
product. Packing is coordinate reinterpretation, so these targets
cost \(O(\kappa)\) extension-field multiplications and
\(O(\kappa^2)\) coordinate operations in total.
The inner-product verifier costs \(O(1+\ell)\) extension-field
operations, including its terminal scalar check.
This gives Eq.~\eqref{eq:pcs-verifier-cost}.

The prover sends \(\kappa\) extension-field elements for
\(\ext{\bm v}\), followed by the \(2\ell+2\) extension-field
elements of Fact~\ref{fact:ip}, including both terminal values.
The verifier computes \(\bm W\) locally, so no matrix message or
base-field proof is sent. Excluding the underlying opening, total
prover-to-verifier communication is therefore
\(\kappa+2\ell+2\) extension-field elements.
\end{proof}

\subsection{Proof of Proposition~\ref{proposition:position-encoding-evaluation}}
\label{appx:proof-position-encoding}

\begin{proof}
Let \(\bm d\in\Field^{2^s}\) list the encoded digits
\(\sum_{t=1}^s\delta_tb_t\) in lexicographic order over
\((b_1,\ldots,b_s)\in\{0,1\}^s\).
Its MLE is:
\begin{equation}
    \left\langle\bm d,\tensorspan{\ext{\bm y}}\right\rangle
       =\sum_{t=1}^{s}\delta_t\ext y_t,
       \qquad \ext{\bm y}\in\ExtensionField^s.
\label{eq:digit-vector-mle}
\end{equation}
The right-hand side is multilinear and agrees with the digit
encoding on the Boolean hypercube, proving the identity over
every characteristic.
Each base-\(\underline q\) digit occupies one consecutive
\(s\)-bit block. Prepending \(\gamma s-\mu\) zero bits restricts
the full digit domain to positions \(0,\ldots,L-1\).
Applying Eq.~\eqref{eq:digit-vector-mle} to each block and taking
the basis-weighted sum proves
Eq.~\eqref{eq:position-encoding-evaluation}.
Only the \(\mu\) original coordinates need to be processed.

For arbitrary tensor factors
\(\ext{\bm\chi}^{(t)}=(\ext\chi^{(t)}_0,\ext\chi^{(t)}_1)^\top\),
the contraction with the \(t\)-th bit function is:
\[
    \ext\chi^{(t)}_1
       \prod_{\substack{u=1\\u\ne t}}^\mu
          (\ext\chi^{(u)}_0+\ext\chi^{(u)}_1).
\]
Prefix and suffix products compute all these terms in
\(O(\mu)\) operations without division. Taking the basis-weighted
sum evaluates the position labels in \(O(\log L)\) operations.
\end{proof}

\subsection{Proof of Theorem~\ref{theorem:single-instance-packing}}
\label{appx:proof-single-instance}

\begin{proof}
\noindent\textbf{Completeness.}
For a satisfying assignment, the five claims in
Eq.~\eqref{eq:plonkWeightedHadamard1} and
Eq.~\eqref{eq:plonkWeightedHadamard2} sum to zero.
The virtual vector \(\bm w_M=\bm w_L\circ\bm w_R\) has its
terminal evaluation checked by a weighted Hadamard claim on the
wire vectors. Completeness of the five Hadamard reductions and
their openings therefore gives completeness of the gate argument.

Linear independence of
\(\ext\theta_1,\ldots,\ext\theta_\gamma,\ext\theta_d\) over
\(\Field\) makes the joint encoding injective in the wire value
and position digits.
In particular, the entry of \(\ext{\bm a}\) with label
\(\ext\tau_i\) can match only the entry of \(\ext{\bm b}\)
at position \(\sigma^{-1}(i)\).
Equality of the two multisets is thus equivalent to
\(w_i=w_{\sigma^{-1}(i)}\) for every \(i\), proving
Eq.~\eqref{eq:together-copy-equivalence}.
Padding adds fixed points and does not change the relation.
Completeness of the permutation and opening protocols completes
the argument.

\smallskip\noindent\textbf{Soundness.}
Condition on binding of the three witness commitments.
Since fixed-basis packing is injective, these commitments determine
three base-field wire vectors and hence the virtual \(\bm w_M\)
before any challenge is sampled.
The selectors and permutation digits are fixed by circuit setup.

If the gate residual is nonzero, its random MLE evaluation
vanishes with probability at most \(\log_2N/|\ExtensionField|\).
Outside this event, an accepting scalar sum contains a false
claim. Its reduction either incurs one of the five Hadamard
errors or leaves a false opening. The same reasoning applies to
the nested \(\bm w_M\) claim: the claimed value is fixed before
its challenges, and the committed wires determine its true value.

If a copy constraint is false, the two encoded multisets differ.
The fingerprint check in Section~\ref{sec:permutation} then has
error at most \(\min\{1,(L-1)/|\ExtensionField|\}\).
Conditioned on a nonzero fingerprint, acceptance requires failure
of one of the two grand-product reductions or a false opening.
Injectivity makes the joint encoding exact.

The gate argument uses five selector and five wire openings.
The two copy terminals add three wire openings each, and the
\(\ext{\bm\tau}_\sigma\) operand adds one digit-table opening
for each of its \(\gamma\) coordinates.
The public-I/O and identity-label evaluations are exact.
Thus there are sixteen length-\(N\) packed openings and
\(\gamma\) length-\(L\) packed openings.
Their adaptive evaluation errors are bounded by
Eq.~\eqref{eq:together-opening-error}.
A union bound over these events and witness-binding failure
gives Eq.~\eqref{eq:together-soundness}.

\smallskip\noindent\textbf{Efficiency.}
The five Hadamard reductions and the sixteen length-\(N\)
packed openings cost \(O(\kappa N)\bb\).
The \(\gamma\) digit-table openings cost
\(O(\gamma\kappa L)\bb\).
Both estimates follow from
Propositions~\ref{proposition:hadamard-complexity} and
\ref{proposition:packed-pcs}, since \(m\geq\kappa\) and
\(L/\kappa\geq m\).

For the copy proofs, a multiplication in \(\mathbb K\) costs
\(O(d^2)\bb\), a \(\mathbb K\)-by-\(\ExtensionField\)
multiplication costs \(O(\kappa d)\bb\), and a multiplication
in \(\ExtensionField\) costs \(O(\kappa^2)\bb\).
Substituting \(L\) and \(\kappa_{\mathbb K}=\kappa/d\)
in Proposition~\ref{proposition:permutation-complexity}
therefore gives:
\[
    O\bigl(\kappa dL+\kappa^2(\log_2L)^2\bigr)\bb.
\]
The length-dependent term absorbs the additive packing costs,
using \(L\geq N\geq\kappa^2\).
Adding the gate, opening, and copy costs proves the prover bound.

For each interpolation check, precomputed polynomial coefficients
allow evaluation of a complete Lagrange basis in
\(O(\kappa^3)\bb\); the tower-field analog is no larger.
The packed openings contribute
\(O((\gamma+1)\kappa^3(1+\log_2L))\bb\)
to verifier time.
The recursive grand-product checks contribute
\(O(\kappa^2(\log_2L)^2)\bb\).
The aligned wire segments require only tensor restrictions,
and the public-I/O contribution can be evaluated in
\(O(n_{\mathsf{io}}+\log L)\) extension-field operations.
Proposition~\ref{proposition:position-encoding-evaluation}
handles the identity labels. These costs give the verifier bound.

The Hadamard and grand-product reductions communicate
\(O(\kappa+(\log_2L)^2)\) extension-field elements.
Adding the packed-opening messages from
Proposition~\ref{proposition:packed-pcs} gives the communication
bound in \(\Field\)-coordinates; the opening term
\(O((\gamma+1)\kappa\log_2L)\) is absorbed by
\(O(\kappa(\log_2L)^2)\), since \(\gamma\leq\log_2L\).
Sequential execution retains \(O(N+Ld+\kappa^2)=O(N+Ld)\)
coordinates, including the digit tables and the
\(\mathbb K\)-valued copy operands.
There are three online witness commitments, and each packed
opening makes one underlying opening call at the stated length.
The theorem adds the underlying PCS costs and workspace separately.
\end{proof}

\subsection{Proof of Theorem~\ref{theorem:multiple-instance-packing}}
\label{appx:proof-multiple-instance}

\begin{proof}
\noindent\textbf{Completeness.}
Let \(\bm R\in\Field^{m\times\kappa}\) contain the gate
residuals of the \(\kappa\) instances as columns.
The left-hand side of Eq.~\eqref{eq:together-multiple-gates}
is exactly \(\bm R\ext{\bm\beta}\), including the term
\(\bm q_C\bm1_\kappa^\top\ext{\bm\beta}\).
Linear independence of the basis implies:
\[
    \bm R\ext{\bm\beta}=\bm0_m
       \quad\Longleftrightarrow\quad
    \bm R=\bm0_{m\times\kappa}.
\]
Similarly, packing commutes with the common row permutation,
so \(\sigma(\ext{\bm w})=\ext{\bm w}\) is equivalent to
the copy constraints in every instance.
This equivalence includes the packed public I/O.
Eq.~\eqref{eq:together-multiple-product-opening} authenticates
the virtual product vector from
Eq.~\eqref{eq:together-multiple-wires}.
All scalar checks and tree relations therefore accept for an
honest prover.

\smallskip\noindent\textbf{Gate soundness.}
Condition on binding of the witness commitments.
Every committed vector in \(\ExtensionField^m\) has a unique
matrix of \(\Field\)-coordinates in the fixed basis.
Thus both the gate residual and \(\ext{\bm w}_M\) are fixed
functions of the committed witnesses.
If some gate is invalid, the packed residual is nonzero.
Its random length-\(m\) tensor projection vanishes with
probability at most \(\log_2m/|\ExtensionField|\).

For a fixed projection point \(\ext{\bm x}\), the gate
sum-check polynomial is the product of the multilinear equality
weight with the residual expression formed from selector and wire
MLEs. Each selector-times-wire term has individual degree two,
and multiplication by the equality weight raises this to three.
The constant-selector term has lower degree, so sum-check has
error at most
\(3\log_2m/|\ExtensionField|\).
Its five selector and three committed-wire evaluations are
checked by direct PCS openings.
The remaining virtual-product evaluation is checked by one
packed Hadamard proof and its two packed openings.
These checks give the first two algebraic terms in
Eq.~\eqref{eq:together-multiple-soundness}. The bound applies
jointly to all \(\kappa\) instances because packing is injective.

\smallskip\noindent\textbf{Copy fingerprint.}
Suppose that the packed assignment violates a copy constraint.
Consider the bivariate difference polynomial:
\[
    P(Y,Z)=
       \prod_{i=1}^{L}(Z+\ext w_i+Y\ext\tau_i)
       -\prod_{i=1}^{L}(Z+\ext w_i+Y\ext\tau_{\sigma,i}).
\]
The labels \(\ext\tau_i\) are distinct because the bit-to-field
digit map is injective and \(\gamma\leq\kappa\) digit slots use
independent basis coordinates.
If \(P\) were zero, unique factorization in
\(\ExtensionField[Y,Z]\) would equate the two multisets of
linear factors, each monic in \(Z\).
Matching their \(Y\)-coefficients would then force
\(\ext w_i=\ext w_{\sigma^{-1}(i)}\) for every \(i\),
contradicting the copy violation. Hence \(P\ne0\).

The homogeneous degree-\(L\) parts are equal:
\[
    \prod_{i=1}^{L}(Z+Y\ext\tau_i)
       =\prod_{i=1}^{L}(Z+Y\ext\tau_{\sigma,i}).
\]
It follows that \(P\) has total degree at most \(L-1\).
The independent uniform challenges \(\ext y,\ext z\) therefore
satisfy \(P(\ext y,\ext z)=0\) with probability at most
\(\min\{1,(L-1)/|\ExtensionField|\}\).
This bound includes the randomness used to combine values and labels.

\smallskip\noindent\textbf{Product-tree soundness.}
For one grand product, unfold the binary tree from its root
to the length-\(L\) input vector.
At a level with \(2^\ell\) parent entries, where
\(\ell\in\{0,\ldots,\mu-1\}\), let
\(\ext{\bm x}\in\ExtensionField^\ell\) be the current
evaluation point.
The weighted product identity uses the polynomial:
\[
    \left\langle
       \tensorspan{\ext{\bm x}},\tensorspan{\bm X}
    \right\rangle
    \left\langle\ext{\bm u},\tensorspan{\bm X}\right\rangle
    \left\langle\ext{\bm v},\tensorspan{\bm X}\right\rangle,
\]
where \(\ext{\bm u},\ext{\bm v}\) are the two child subvectors
and \(\bm X\) denotes \(\ell\) formal variables.
Its degree in each variable is at most three.
Sum-check reduces its Boolean sum to child evaluations
\(\ext v_L,\ext v_R\) at one point \(\ext{\bm t}\), with
error at most \(3\ell/|\ExtensionField|\).

After both values are fixed, sample a fresh
\(\ext\eta\sample\ExtensionField\) and form the next claim:
\[
    (1-\ext\eta)\ext v_L+\ext\eta\ext v_R
       =
    \left\langle
       \ext{\bm u}\mathbin{\|}\ext{\bm v},
       \tensorspan{(\ext\eta,\ext{\bm t})}
    \right\rangle.
\]
If either child claim is false, this merge accepts it
with probability at most \(1/|\ExtensionField|\).
Thus one tree has reduction error at most:
\[
    \sum_{\ell=0}^{\mu-1}
       \frac{3\ell+1}{|\ExtensionField|}
       =
    \frac{3\mu(\mu-1)/2+\mu}{|\ExtensionField|}.
\]
The \(\ell=0\) level has no sum-check rounds but still uses
a fresh merge challenge.
Each tree ends in a single MLE claim on its original input.
Subtracting the known shift \(\ext z\) and expanding the
value-label combination produces three wire claims per tree
and one permuted-label claim across the two trees.

\smallskip\noindent\textbf{Openings and composition.}
The gate argument uses eight direct openings and two packed
openings. The copy argument uses seven direct openings.
Fourteen direct openings concern length-\(m\) vectors and
one concerns \(\ext{\bm\tau}_\sigma\) of length \(L\).
Each packed opening terminates in one underlying opening at
length \(m\). Consequently the total opening error is at most
\(15\epsilon_{\mathsf{PCS}}+2\epsilon_{\mathsf{open}}(N)\).
The PCS assumption covers these adaptive queries.
A union bound with the gate, fingerprint, two tree, and binding
errors proves Eq.~\eqref{eq:together-multiple-soundness}.

\smallskip\noindent\textbf{Efficiency.}
Computing \(\bm W_L\circ\bm W_R\) costs \(O(N)\bb\).
The one packed Hadamard proof and two packed openings cost
\(O(\kappa N)\bb\), using \(m\geq\kappa\).
The gate sum-check operates on a constant number of
length-\(m\) extension-field tables and costs
\(O(m)\) extension-field operations.

For the product trees, the layer lengths form a geometric
series summing to \(O(L)\). Forming their entries, generating
the equality tables, and running the layer sum-checks therefore
cost \(O(L+(\log_2L)^2)\) extension-field operations in total.
The layer transcripts and verifier checks cost
\(O((\log_2L)^2)\) extension-field elements and operations,
respectively.
Schoolbook multiplication converts these bounds to the
\(\kappa^2[L+(\log_2L)^2]\bb\) term.

The packed Hadamard check contributes \(O(\kappa^3)\bb\)
to verifier time, and the two packed opening reductions add
\(O(\kappa^3(1+\log_2m))\bb\).
Packing the public I/O reads \(O(\kappa n_{\mathsf{io}})\)
base-field coordinates, and its terminal contractions cost
\(O(n_{\mathsf{io}}+\log L)\) extension-field operations.
The identity labels are handled by
Proposition~\ref{proposition:position-encoding-evaluation}.
Combining these costs proves the verifier bound.

The packed messages occupy \(O(\kappa^2+\kappa\log_2m)\)
base-field coordinates, in addition to the
\(O((\log_2L)^2)\) extension-field elements of the direct
gate and copy reductions.
Sequential execution stores \(O(N)\) witness coordinates and
\(O(L)\) extension-field product-tree entries, giving
\(O(N+\kappa L)\) coordinates of space.
These are the bounds in
Eq.~\eqref{eq:together-multiple-efficiency}.
The three witness commitments and seventeen underlying openings
are charged separately, at the lengths stated in the theorem.
For \(L=O(m)\), substituting \(N=m\kappa\) and dividing the
total prover cost by \(\kappa\) gives the amortized bound.
\end{proof}
